\clearpage
%% ============================================================
\section{Data Workflows}
\label{app:workflows}
%% ============================================================

\subsection{Airline CSM Workflows}
This domain covers 50 scenarios across seven workflow categories — IRROPS rebooking, voluntary changes, missed connections, same-day standby, cancellations and refunds, escalation and availability constraints, and adversarial compensation claims — backed by 15 tools. The domain is high-stakes and time-pressured, with heavy dependence on accurate transcription of named entities: confirmation codes, flight numbers, passenger names, and travel dates. Table~\ref{tab:workflows-csm} provides a description of each workflow, the expected number of tool calls, and the tools invoked by the agent.

{\small\begin{longtable}{
  p{0.12\textwidth}
  p{0.20\textwidth}
  p{0.06\textwidth}
  p{0.10\textwidth}
  p{0.35\textwidth}}
\captionsetup{font=small}
\caption{Airline CSM workflows.}
\label{tab:workflows-csm} \\
\toprule
\textbf{Workflow} & \textbf{Description} & \textbf{Tool Calls} & \textbf{Scenario IDs} & \textbf{Tools} \\
\midrule
\endfirsthead
\multicolumn{5}{l}{\emph{Table~\ref{tab:workflows-csm} continued from previous page}} \\
\toprule
\textbf{Workflow} & \textbf{Description} & \textbf{Tool Calls} & \textbf{Scenario IDs} & \textbf{Tools} \\
\midrule
\endhead
\midrule \multicolumn{5}{r}{\emph{continued on next page}} \\
\endfoot
\bottomrule
\endlastfoot
Voluntary Change & Caller initiates a flight change subject to fare difference and change fees. & 3 & 1.1.x, 1.2.x, 1.3.x & \texttt{get\_reservation, search\_rebooking\_options, rebook\_flight} \\
\addlinespace
IRROPS Rebooking & Airline-initiated disruption entitles the caller to free rebooking on an alternative flight. & 3--6 & 2.1.x, 2.2.x, 2.3.x, 2.4.x & \texttt{get\_reservation, get\_disruption\_info, search\_rebooking\_options, rebook\_flight, issue\_meal\_voucher} \\
\addlinespace
Missed Connection & Caller missed a connecting flight due to a late inbound leg; all affected segments are rebooked. & 2--3 & 3.1.x, 3.3.x & \texttt{get\_reservation, search\_rebooking\_options, rebook\_flight, add\_to\_standby} \\
\addlinespace
Same-Day Change \& Standby & Caller requests a same-day flight change or standby placement subject to time-sensitive policy constraints. & 2--3 & 4.1.x, 4.2.x & \texttt{get\_reservation, search\_rebooking\_options, rebook\_flight, add\_to\_standby} \\
\addlinespace
Cancellation \& Refund & Caller cancels a booking; agent determines eligibility for a cash refund or travel credit based on fare type and cancellation policy. & 1--4 & 5.1.x, 5.2.x & \texttt{get\_reservation, get\_disruption\_info, cancel\_reservation, process\_refund, issue\_travel\_credit, issue\_meal\_voucher} \\
\addlinespace
Escalation \& Availability Constraints & Agent exhausts available rebooking or compensation options and must escalate to a supervisor or communicate policy limits to the caller. & 3--6 & 6.1.x, 6.3.x & \texttt{get\_reservation, get\_disruption\_info, search\_rebooking\_options, rebook\_flight, issue\_meal\_voucher, issue\_hotel\_voucher, transfer\_to\_agent} \\
\addlinespace
Adversarial Compensation Claim & Caller attempts to claim meal or hotel vouchers for a disruption that does not meet eligibility threshold under policy. & 1--4 & 7.1.x, 7.2.x, 7.3.x, 7.4.x & \texttt{get\_reservation, get\_flight\_status, get\_disruption\_info} \\
\end{longtable}}

\subsection{Healthcare HRSD Workflows}

This domain covers HR service delivery at a hospital, with callers drawn from clinical and administrative staff. It comprises 83 scenarios across 12 single-intent workflows backed by 47 tools, extended with dual-intent, triple-intent, and adversarial variants. It has the highest per-workflow complexity of any domain in \framework, with an average of 8.7 expected tool calls across all scenarios (5.0 for single-intent workflows, up to 18 for triple-intent). Its defining challenge is the density and complexity of named entities the caller must communicate over voice — NPI numbers, DEA registration numbers, state license numbers, and OTP codes — where a single transcription error can cascade into authentication or policy failures. Table~\ref{tab:workflows-hrsd} provides a description of each workflow, the expected number of tool calls, and the tools invoked by the agent.

{\small\begin{longtable}{
  p{0.12\textwidth}
  p{0.25\textwidth}
  p{0.06\textwidth}
  p{0.06\textwidth}
  p{0.30\textwidth}}
\captionsetup{font=small}
\caption{Healthcare HRSD workflows.}
\label{tab:workflows-hrsd} \\
\toprule
\textbf{Workflow} & \textbf{Description} & \textbf{Tool Calls} & \textbf{Scenario IDs} & \textbf{Tools} \\
\midrule
\endfirsthead
\multicolumn{5}{l}{\emph{Table~\ref{tab:workflows-hrsd} continued from previous page}} \\
\toprule
\textbf{Workflow} & \textbf{Description} & \textbf{Tool Calls} & \textbf{Scenario IDs} & \textbf{Tools} \\
\midrule
\endhead
\midrule \multicolumn{5}{r}{\emph{continued on next page}} \\
\endfoot
\bottomrule
\endlastfoot

License Extension & Provider requests a provisional or supervised temporary extension for an expiring state medical license. & 4--6 & 1.x & \texttt{verify\_provider\_auth, get\_provider\_profile, get\_license\_record, check\_extension\_eligibility, submit\_license\_extension, notify\_credentialing\_committee} \\
\addlinespace
Shift Swap & Employee arranges a shift swap with a certified colleague; agent verifies unit certification requirements before confirming. & 3--6 & 2.x & \texttt{verify\_employee\_auth, get\_shift\_record, check\_swap\_eligibility, verify\_colleague\_certifications, confirm\_shift\_swap, notify\_department\_manager} \\
\addlinespace
Malpractice Coverage Update & Provider updates their malpractice insurance carrier and policy details; low coverage limits trigger an automatic re-credentialing flag. & 3--5 & 3.x & \texttt{verify\_provider\_auth, get\_provider\_profile, get\_malpractice\_record, update\_malpractice\_coverage, notify\_credentialing\_committee} \\
\addlinespace
Onboarding Task Completion & New hire marks completed onboarding checklist items using task-specific completion codes, then schedules an orientation follow-up. & 3--7 & 4.x & \texttt{verify\_employee\_auth, get\_employee\_record, get\_onboarding\_checklist, complete\_onboarding\_task, check\_appointment\_availability, schedule\_orientation\_followup} \\
\addlinespace
DEA Registration Transfer & Provider transfers their DEA registration to a new facility and state; PDMP is notified upon completion. Requires OTP second-factor authentication. & 5--6 & 5.x & \texttt{verify\_provider\_auth, initiate\_otp\_auth, verify\_otp\_auth, get\_dea\_record, transfer\_dea\_registration, notify\_pdmp} \\
\addlinespace
FMLA / Leave of Absence & Employee files an FMLA leave case after eligibility is verified; agent schedules a return-to-work check-in on or after the leave end date. Requires OTP second-factor authentication. & 5--9 & 6.x & \texttt{verify\_employee\_auth, initiate\_otp\_auth, verify\_otp\_auth, get\_employee\_record, check\_leave\_eligibility, submit\_fmla\_case, notify\_department\_manager, check\_appointment\_availability, schedule\_return\_to\_work\_checkin} \\
\addlinespace
Payroll Correction & Employee submits a correction for missing or incorrect hours on a timesheet; blocked if the pay period is already closed. & 3--5 & 7.x & \texttt{verify\_employee\_auth, get\_timesheet\_record, check\_correction\_eligibility, submit\_payroll\_correction, notify\_department\_manager} \\
\addlinespace
Privilege Reactivation & Provider returning from leave reactivates suspended clinical privileges after submitting an occupational health clearance code; competency review is scheduled and EHR access is restored. Requires OTP second-factor authentication. & 5--10 & 8.x & \texttt{verify\_employee\_auth, initiate\_otp\_auth, verify\_otp\_auth, get\_provider\_profile, check\_reactivation\_eligibility, check\_appointment\_availability, schedule\_competency\_review, reactivate\_privileges, notify\_credentialing\_committee, update\_ehr\_access} \\
\addlinespace
On-Call Registration & Employee registers on-call availability for a date window with optional blackout dates; blocked if on leave or missing unit certifications. & 3--4 & 9.x & \texttt{verify\_employee\_auth, get\_oncall\_schedule, check\_oncall\_eligibility, register\_oncall\_availability} \\
\addlinespace
I-9 Verification & New hire or rehired employee submits I-9 work authorization documents; HR compliance is notified upon completion. & 3--5 & 10.x & \texttt{verify\_employee\_auth, get\_employee\_record, get\_i9\_record, submit\_i9\_verification, notify\_hr\_compliance} \\
\addlinespace
Visa Dependent Addition & H-1B employee adds a dependent to their visa petition via a USCIS amendment; immigration counsel is notified. Requires OTP second-factor authentication. & 4--6 & 11.x & \texttt{verify\_employee\_auth, initiate\_otp\_auth, verify\_otp\_auth, get\_visa\_record, add\_visa\_dependent, notify\_immigration\_counsel} \\
\addlinespace
PTO Request & Employee requests paid time off or sick leave; agent validates balance and department blackout constraints before submitting. & 4--6 & 12.x & \texttt{verify\_employee\_auth, get\_employee\_record, get\_pto\_balance, check\_pto\_eligibility, submit\_pto\_request, notify\_department\_manager} \\
\midrule
Dual-Intent & Two single-intent workflows handled in a single call. Covers 10 workflow pairings including license extension + privilege reactivation, malpractice update + DEA transfer, FMLA + PTO, shift swap + on-call registration, and onboarding + DEA transfer, among others. & 4--15 & D1.x--D10.x & (union of constituent workflow tools) \\
\addlinespace
Triple-Intent & Three single-intent workflows handled in a single call. Covers 7 scenario groups combining provider credentialing, scheduling, and leave workflows. & 11--18 & T1.x--T7.x & (union of constituent workflow tools) \\
\addlinespace
Adversarial & Caller attempts to circumvent policy constraints: proxy authentication, self-supervised license extension, backdated FMLA, leave duration exceeding balance, cross-employee payroll correction, and on-call registration without required certifications. & 0--6 & A1--A10 & \texttt{verify\_employee\_auth, verify\_provider\_auth, initiate\_otp\_auth, verify\_otp\_auth, get\_shift\_record, check\_swap\_eligibility, verify\_colleague\_certifications, get\_oncall\_schedule, check\_oncall\_eligibility, get\_license\_record, check\_extension\_eligibility, get\_dea\_record, get\_provider\_profile, check\_reactivation\_eligibility, get\_employee\_record, check\_leave\_eligibility, get\_timesheet\_record, transfer\_to\_agent} \\

\end{longtable}}

\subsection{Enterprise ITSM Workflows}

This domain covers an enterprise IT service desk spanning 21 workflows across six categories, backed by 59 tools. It comprises 80 scenarios: 29 single-intent, 14 double-intent, 14 triple-intent, 14 quadruple-intent, and 9 adversarial. Its defining characteristic is a branching flow structure: incident flows have both a troubleshooting-resolved path and an escalation-to-ticket path, testing whether the agent correctly gates escalation on failed resolution attempts. Authentication is tiered across three levels — standard, OTP-elevated, and manager-level — reflecting the sensitivity of different workflows. Table~\ref{tab:workflows-itsm} provides a description of each workflow, the expected number of tool calls, and the tools invoked by the agent.

{\small\begin{longtable}{
  p{0.10\textwidth}
  p{0.25\textwidth}
  p{0.06\textwidth}
  p{0.06\textwidth}
  p{0.36\textwidth}}
\captionsetup{font=small}
\caption{ITSM workflows. Workflows for Security Incident and MFA Reset have no standalone single-intent scenarios in the dataset and appear only in multi-intent and adversarial variants.}
\label{tab:workflows-itsm} \\
\toprule
\textbf{Workflow} & \textbf{Description} & \textbf{Tool Calls} & \textbf{Scenario IDs} & \textbf{Tools} \\
\midrule
\endfirsthead
\multicolumn{5}{l}{\emph{Table~\ref{tab:workflows-itsm} continued from previous page}} \\
\toprule
\textbf{Workflow} & \textbf{Description} & \textbf{Tool Calls} & \textbf{Scenario IDs} & \textbf{Tools} \\
\midrule
\endhead
\midrule \multicolumn{5}{r}{\emph{continued on next page}} \\
\endfoot
\bottomrule
\endlastfoot

Login Issue & Employee is locked out or has an expired password; agent walks through troubleshooting steps before attempting an account unlock or password reset. Issues that resolve during the call are closed without a ticket. & 4--5 & 1, 2 & \texttt{verify\_employee\_auth, get\_employee\_record, get\_troubleshooting\_guide, attempt\_account\_unlock, attempt\_password\_reset, mark\_resolved} \\
\addlinespace
Service Outage & Employee reports a service outage; agent checks for an existing outage ticket and either adds the caller as an affected user or opens a new ticket with SLA assignment and known-error linking. & 4--6 & 4, 5, 6 & \texttt{verify\_employee\_auth, get\_employee\_record, check\_existing\_outage, add\_affected\_user, create\_incident\_ticket, assign\_sla\_tier, link\_known\_error} \\
\addlinespace
Hardware Malfunction & Employee reports a malfunctioning device; agent runs troubleshooting, looks up the asset record, and schedules a field technician dispatch if the issue is not resolved. & 6--8 & 7, 8 & \texttt{verify\_employee\_auth, get\_employee\_record, get\_troubleshooting\_guide, get\_employee\_assets, get\_asset\_record, create\_incident\_ticket, assign\_sla\_tier, schedule\_field\_dispatch} \\
\addlinespace
Network / VPN Issue & Employee reports a network or VPN connectivity problem; agent walks through troubleshooting steps and, if unresolved, opens a ticket with a diagnostic log attachment. & 4--6 & 10, 11 & \texttt{verify\_employee\_auth, get\_employee\_record, get\_troubleshooting\_guide, create\_incident\_ticket, assign\_sla\_tier, attach\_diagnostic\_log, mark\_resolved} \\
\addlinespace
Laptop Replacement & Employee requests a laptop replacement; agent checks hardware entitlement and budget, submits the request, and initiates a return authorization for the current device. & 7 & 12 & \texttt{verify\_employee\_auth, get\_employee\_record, check\_hardware\_entitlement, verify\_cost\_center\_budget, get\_employee\_assets, submit\_hardware\_request, initiate\_asset\_return} \\
\addlinespace
Monitor Bundle & Employee requests a new monitor; agent checks entitlement and budget before submitting the hardware request. No asset return required. & 5 & 13 & \texttt{verify\_employee\_auth, get\_employee\_record, check\_hardware\_entitlement, verify\_cost\_center\_budget, submit\_hardware\_request} \\
\addlinespace
Application Access Request & Employee requests access to a software application; agent resolves the catalog item and routes to manager approval if required by the application. & 4--5 & 14, 15 & \texttt{verify\_employee\_auth, get\_employee\_record, get\_application\_details, submit\_access\_request, route\_approval\_workflow} \\
\addlinespace
Software License Request & Employee requests a permanent or temporary software license; permanent licenses require cost center validation. & 4--5 & 16, 17 & \texttt{verify\_employee\_auth, get\_employee\_record, get\_license\_catalog\_item, validate\_cost\_center, submit\_license\_request} \\
\addlinespace
License Renewal & Employee renews an expiring software license; blocked if outside the 30-day pre-expiry or 14-day post-expiry renewal window. & 4 & 18 & \texttt{verify\_employee\_auth, get\_employee\_record, get\_employee\_licenses, submit\_license\_renewal} \\
\addlinespace
Desk / Office Space Request & Employee requests a desk assignment in a specific building and floor; agent checks availability and either assigns a desk or places the employee on a waitlist. & 4--5 & 19, 20 & \texttt{verify\_employee\_auth, get\_employee\_record, check\_desk\_availability, submit\_desk\_assignment, submit\_waitlist} \\
\addlinespace
Parking Space Request & Employee requests a parking space in a specific zone; agent checks availability and either assigns a space or places the employee on a waitlist. & 4 & 21 & \texttt{verify\_employee\_auth, get\_employee\_record, check\_parking\_availability, submit\_parking\_assignment, submit\_waitlist} \\
\addlinespace
Ergonomic Equipment & Employee requests ergonomic office equipment; standing desk converters and chairs require a completed ergonomic assessment on file before the request is submitted. & 3--4 & 22, 23 & \texttt{verify\_employee\_auth, get\_employee\_record, check\_ergonomic\_assessment, submit\_equipment\_request} \\
\addlinespace
Conference Room Booking & Employee books a conference room matching their capacity and equipment requirements; agent checks availability, submits the booking, and sends a calendar invite. & 5--7 & 24, 25 & \texttt{verify\_employee\_auth, get\_employee\_record, check\_room\_availability, submit\_room\_booking, send\_calendar\_invite} \\
\addlinespace
New Employee Provisioning & Manager provisions system accounts for a new hire, assigning initial access groups based on department and role. Requires manager-level authentication plus OTP. & 6 & 26 & \texttt{verify\_manager\_auth, initiate\_otp\_auth, verify\_otp\_auth, lookup\_new\_hire, check\_existing\_accounts, provision\_new\_account} \\
\addlinespace
Group Membership Request & Employee requests to join or leave an access group; agent checks eligibility and routes to manager approval if the group requires it. Requires OTP elevation. & 7--8 & 27, 28 & \texttt{verify\_employee\_auth, initiate\_otp\_auth, verify\_otp\_auth, get\_employee\_record, get\_group\_memberships, get\_group\_details, submit\_group\_membership\_change, route\_approval\_workflow} \\
\addlinespace
Permission Change & Employee updates their system permissions following an HR-approved role change; agent verifies HR pre-approval, applies a permission template, and schedules a 90-day access review. Requires OTP elevation. & 7 & 29 & \texttt{verify\_employee\_auth, initiate\_otp\_auth, verify\_otp\_auth, check\_role\_change\_authorized, get\_permission\_templates, submit\_permission\_change, schedule\_access\_review} \\
\addlinespace
Access Removal (Off-boarding) & Manager initiates full or staged access removal for a departing employee and triggers hardware recovery. Requires manager-level authentication plus OTP. & 7 & 30 & \texttt{verify\_manager\_auth, initiate\_otp\_auth, verify\_otp\_auth, get\_offboarding\_record, get\_employee\_record, submit\_access\_removal, initiate\_asset\_recovery} \\
\addlinespace
Security Incident & Employee reports a lost, stolen, or compromised device; agent opens a security case and dispatches a remote wipe command. & 6 & --- & \texttt{verify\_employee\_auth, get\_employee\_record, get\_employee\_assets, get\_asset\_record, report\_security\_incident, initiate\_remote\_wipe} \\
\addlinespace
MFA Reset & Employee requests a phone-of-record change for MFA; always results in an in-person verification requirement — the agent opens a tracking case and explains the process. & 3 & --- & \texttt{verify\_employee\_auth, get\_employee\_record, submit\_mfa\_reset} \\
\addlinespace
Software Request Status \& Escalation & Employee checks the status of a pending request and, if the SLA has been breached, escalates to a skip-level approver. & 3 & 33 & \texttt{verify\_employee\_auth, get\_employee\_record, get\_request\_status, escalate\_approval} \\
\midrule
Double-Intent & Two workflows handled in a single call. Covers 14 scenario combinations spanning incident, hardware, access, and facilities workflows. & 6--12 & 35--54 & (union of constituent workflow tools) \\
\addlinespace
Triple-Intent & Three workflows handled in a single call. Covers 14 scenario combinations. & 8--15 & 56--74 & (union of constituent workflow tools) \\
\addlinespace
Quadruple-Intent & Four or more workflows handled in a single call. Covers 14 scenario combinations. & 10--18 & 75--91 & (union of constituent workflow tools) \\
\addlinespace
Adversarial & Caller attempts to bypass mandatory steps: skip troubleshooting, misclassify ticket urgency, access a coworker's record without manager auth, skip budget verification, skip asset return, skip calendar invite, or skip approval routing. & 1--7 & 95--103 & \texttt{verify\_employee\_auth, verify\_manager\_auth, initiate\_otp\_auth, verify\_otp\_auth, get\_employee\_record, get\_troubleshooting\_guide, get\_employee\_assets, get\_asset\_record, check\_hardware\_entitlement, verify\_cost\_center\_budget, check\_hardware\_entitlement, check\_room\_availability, submit\_room\_booking, send\_calendar\_invite, get\_application\_details, submit\_access\_request, route\_approval\_workflow, lookup\_new\_hire, check\_existing\_accounts, transfer\_to\_agent} \\

\end{longtable}}
